\section{Our Approach}
\label{sec:impl}

We consider five configurations that grow the same exact CART or C4.5
tree. They differ in how a node's sorted order is maintained, how
candidate cuts are evaluated, and where nodes are processed.
\ImplGpu{} and \ImplHybrid{} are two configurations of one CUDA
backend, without and with CPU hand-offs.

\subsection{Common Design}
\label{sec:impl-common}

The data are stored column-major as single-precision values, with
class labels mapped to small integer identifiers. All configurations
score cuts with the same code, compiled for CPU and GPU alike:
weighted Gini impurity for CART, weighted entropy for C4.5. Gains
differing by at most $10^{-12}$ are equal; the earlier cut, then the
lower attribute index, prevails. Optimised implementations discard
only cuts that cannot replace the incumbent under this rule;
surviving candidates are scored exactly. All configurations
therefore grow identical trees, as checked against the verified
trees of Sect.~\ref{sec:intro}.

Except for \ImplSerialBase{}, an attribute is represented by 8-byte
entries pairing a value with its sample and class packed into one
word, and a node occupies the same index interval in every column.
Once the columns are sorted, a stable partition of this interval
yields the two sorted child intervals, so no attribute is sorted
again---the presorting principle of SLIQ and
SPRINT~\cite{mehta1996sliq,shafer1996sprint}, with the class in the
entry so that sweeps read sequentially. The tree is stored in flat
preorder arrays; concurrently created nodes are held in a shared
store and renumbered at the end.

\subsection{\ImplSerialBase}
\label{sec:impl-simple}

\ImplSerialBase{} is the textbook algorithm on one thread. At each
node it copies and sorts the node's sample identifiers for every
attribute, evaluates every admissible cut exactly, and applies the
shared CART or C4.5 choice rule; the samples are then divided into
two new identifier lists and the children grown recursively. Sorting
a node of $n$ samples costs $O(F n \log n)$ for $F$ attributes, with
values read at random positions through the identifiers, so it is
bound by memory latency rather than bandwidth.

\subsection{\ImplSerialOpt}
\label{sec:impl-serial}

\ImplSerialOpt{} changes the execution, not the algorithm, in three
ways. \emph{Sort once.} Each attribute is ordered once by a stable
radix sort---on one thread in this configuration---and a split is
then made by position: the first $leftCount$ entries of the winning
column form the left child, the corresponding samples are marked,
and every other column is partitioned stably into its child
intervals without comparing values against the threshold again
(large nodes out of place into a spare slot, small nodes in place).
Attributes constant over a node are dropped for its whole subtree.

\emph{Score few cuts exactly.} The sweep maintains class counts
incrementally and first estimates the gain cheaply. Only cuts that
can still beat the best exact gain under the tie rule are rescored
exactly, in cut order; for two classes, blocks whose corner gains
cannot win are skipped. \emph{Allocate once.} The columns occupy
$F+1$ slots for the whole run, an explicit stack replaces recursion,
and large allocations request transparent huge pages. This is the
builder of \ImplParallel{} run without a thread pool.

\subsection{\ImplParallel}
\label{sec:impl-parallel}

\ImplParallel{} adds a thread pool to this builder; layout, sweep
and partition are unchanged, and, as in the NAP strategy of Aldinucci
et al.~\cite{aldinucci2014fastflow}, the kind of parallelism is
chosen by node size. For nodes of at least 32{,}768 samples, the
attributes are swept in parallel and the columns partitioned
concurrently, one per thread; below this size, waking the pool would
cost more than the work. After a node of at least 4{,}096 samples
has been split, one child becomes a pool task while the current
thread continues with the other, and the grower waits only once, at
the end. Presort is parallel too, and spare slots grow with the
thread count, up to $F-1$.

We make no explicit NUMA placement for CPU growth: workers are not
pinned, columns are not divided by socket, and page placement is
left to the operating system. The thread count follows the CPUs the
process may use, so placement can be imposed from outside, e.g.\
with \texttt{taskset} or \texttt{numactl}.

\subsection{\ImplGpu}
\label{sec:impl-gpu}

\ImplGpu{} is the CUDA backend with the hand-off threshold set to
one sample (\texttt{--cuda --gpu-only}), so no subtree is delegated
to the CPU. Whereas the CPU builders grow one node at a time,
depth-first, the GPU processes all nodes of one tree level together.
A level's work is divided into tiles of 2{,}048 consecutive entries
of one node and attribute, one thread block per tile; a fixed
sequence of kernels detects constant segments, sweeps the tiles for
the best cut of every attribute, and partitions the columns for the
children. The
best cut of each attribute is returned to the host, which applies
the same CART or C4.5 choice rule as on the CPU and finishes leaf
children: the only host round trip of a level. The stable partition
marks only the smaller side of each split in a one-bit-per-sample
bitmap.

On the data-center GPUs we evaluate, every cut is scored in double
precision by the shared routine. For GPUs with slow double
precision, an optional single-precision pass discards cuts that
cannot win under a conservative error bound; only survivors are
rescored exactly. Uploads use page-locked staging buffers, allocated
on the GPU's NUMA node (identified through its PCI device), with the
copying threads pinned to that node's CPUs and the transfers
overlapped with the device presort. This is the only explicit NUMA
placement in our implementations, serving both GPU configurations.
The CPU stages transfers and runs post-processing, but grows no
part of the tree.

\subsection{\ImplHybrid}
\label{sec:impl-hybrid}

\ImplHybrid{} (\texttt{--cuda}) uses the same level-wise growth, but
children below a size threshold are finished by the CPU pool while
the GPU continues with the next level
(Fig.~\ref{fig:approach-hybrid}). The threshold can be set
explicitly; by default it grows slowly with the training-set size,
as $N/8192$ clamped to $[1024, 2048]$ samples for CART and $N/4096$
clamped to $[1024, 4096]$ for C4.5. Size is the only criterion: a
GPU level costs about as much for many small nodes as for a few
large ones, whereas a small subtree is an independent CPU task.

\begin{figure}[t]
  \centering
  \includegraphics[width=\textwidth]{figs/approach-hybrid.pdf}
  \caption{One level of \ImplHybrid{}. The GPU sweeps and partitions
  all large nodes of a level together; children at or above the
  threshold $m$ form the next GPU level. Children below $m$ are
  gathered into a device ring, copied through page-locked buffers
  into a host arena, and grown by the CPU pool while the GPU
  proceeds with the next level. Both devices write nodes of the same
  tree to a shared store, renumbered in preorder at the end.
  \ImplGpu{} is the same scheme with $m=1$.}
  \label{fig:approach-hybrid}
\end{figure}

Handed-off children are gathered on the device, restricted to
attributes that can still split them, into a ring buffer, and copied
on separate streams through page-locked buffers into a pageable host
arena. A host callback queued behind each copy starts the CPU tasks
once their data have arrived, so the level loop never waits for a
hand-off. The arena is not bound to the GPU's NUMA node.
\ImplHybrid{} thus neither assigns whole trees of a forest to a
device nor reduces the CPU to a controller: individual nodes of one
tree go to the device that grows them at the lower marginal cost.
